\subsection{Implementierung der Raum- und Belegungsverwaltung (5123097)}

Die Raumverwaltung stellt eine zentrale Komponente des Hospital Management Systems dar. 
Sie umfasst die hierarchische Struktur von Abteilungen (\textit{Departments}), Stationen (\textit{Stations}) und Zimmern (\textit{Rooms}) sowie die darauf aufbauenden Belegungen (\textit{Bookings}). 
Bei der Umsetzung wurde darauf geachtet, die Netzwerklast zu optimieren und gleichzeitig eine hohe Datenkonsistenz zu gewährleisten.

\subsubsection{Interne Objektverknüpfung vs. API-Effizienz}
Innerhalb des Backends werden die Entitäten als vollständig verknüpfte Objektgraphen verwaltet. 
Dies ermöglicht es der Geschäftslogik, effizient auf Relationen zuzugreifen (z. B. von einer Buchung direkt auf das Zimmer-Objekt und dessen Kapazitäten zuzugreifen), ohne manuelle Datenbank-Joins im Code formulieren zu müssen.

%TODO evtl im dritten Portfolio mit einbeziehen falls das dann umgesetzt
%Um jedoch den \textbf{Network Traffic} zwischen Backend und Frontend zu minimieren, folgt die API einer bedarfsorientierten Strategie:
%\begin{itemize}
%    \item \textbf{ID-Referenzierung:} In der API-Antwort werden tieferliegende oder für den aktuellen Kontext nicht unmittelbar relevante Relationen primär als IDs übertragen. 
%    Dies verhindert das Übertragen massiver Datenmengen (z. B. vollständige Abteilungsdetails bei jeder einzelnen Zimmerabfrage), wenn der Client diese Informationen bereits lokal hält oder nicht benötigt.
%    \item \textbf{Gezielte Payload-Reduktion:} Durch den Einsatz dedizierter Transferstrukturen wird sichergestellt, dass nur die für die jeweilige Ansicht essenziellen Datenfelder serialisiert werden. 
%    Dies führt zu einer performanten Schnittstelle, die besonders bei instabilen Netzwerkverbindungen im Krankenhausalltag (z. B. auf mobilen Visiten-Tablets) zuverlässig arbeitet.
%\end{itemize}



\subsubsection{Request-Objekte bei Schreibzugriffen}
Im Gegensatz dazu werden für Schreiboperationen (\texttt{POST}) dedizierte \textbf{Request-Objects} (z. B. \texttt{BookingRequest}) verwendet. 
Diese enthalten zur Reduktion der Komplexität und Payload-Größe lediglich die Primärschlüssel (\textit{IDs}) der verknüpften Entitäten. 
Dies entkoppelt die API von der internen Datenbankstruktur und verhindert, dass Clients fälschlicherweise versuchen, Sub-Ressourcen wie Abteilungsnamen über einen \texttt{Bookings}-Endpunkt zu ändern.

\subsubsection{Endpoints}
Die folgenden Beispiele illustrieren die Verwendung der Endpunkte für Abfragen und Ressourcen-Erstellung:

\begin{itemize}
    \item \textbf{Abteilungen und Stationen abrufen:}
    \begin{itemize}
        \item \texttt{GET /api/departments?page=0\&size=10\&building=4B}: Abruf von Abteilungen in einem spezifischen Gebäude mit Paginierung.
        \item \texttt{GET /api/stations?nameContains=Intensiv\&departmentId=1}: Gefilterte Suche nach Stationen innerhalb einer Abteilung.
    \end{itemize}
    
    \item \textbf{Raum- und Belegungsinformationen:}
    \begin{itemize}
        \item \texttt{GET /api/rooms/1}: Abruf eines spezifischen Zimmers basierend auf dessen ID.
        \item \texttt{GET /api/bookings?state=COMPLETED\&page=2}: Abruf abgeschlossener Buchungen über die Paginierungsschnittstelle.
    \end{itemize}
    
    \item \textbf{Erstellung und Verwaltung (Schreibzugriffe):}
    \begin{itemize}
        \item \texttt{POST /api/bookings}: Erstellung einer neuen Belegung unter Verwendung eines \textit{BookingRequest}-Objekts (enthält \texttt{room\_id} und \texttt{patient\_id}).
    \end{itemize}
\end{itemize}


\textbf{Beispiel für eine Response (\texttt{\detokenize{GET /api/stations/{STATION_ID}}}):}

\begin{lstlisting}
{
  "(*\color{keyblue}id*)": (*\color{numorange}1*),
  "(*\color{keyblue}name*)": "(*\color{valgreen}Interdisziplinaere Intensivstation*)",
  "(*\color{keyblue}department*)": {
    "(*\color{keyblue}id*)": (*\color{numorange}1*),
    "(*\color{keyblue}name*)": "(*\color{valgreen}Klinik fuer Anaesthesiologie*)",
    "(*\color{keyblue}building*)": "(*\color{valgreen}4B*)"
  }
}
\end{lstlisting}

Im Zuge der Architekturentwicklung wurde die Struktur der Endpunkte kritisch hinterfragt.\\
In einer ersten Version befanden sich die Funktionen zur Verlegung und Entlassung eines Patienten innerhalb des \texttt{bookings}-Endpunkts (u.a. \texttt{POST /api/bookings/relocate?patientId=\detokenize{PATIENT_ID}}). \\
Um die REST-Architektur jedoch möglichst logisch und konsistent aufzubauen, wurde dieser Use Case in den \texttt{patient}-Endpunkt verschoben, da eine Verlegung bzw. Entlassung primär ein Ereignis im Lebenszyklus eines Patienten ist. 
Durch diese Verschiebung wird der \texttt{bookings}-Endpunkt auf die reine Verwaltung von Raumreservierungen reduziert, während patientenzentrierte Aktionen konsistent an der Patienten-Ressource gebündelt werden. 
%Dies verbessert die Intuitivität der API für Entwickler und folgt dem Prinzip der \textit{Separation of Concerns}.